% Chapter: the approximate inverse for the horizontal problem on T^2 and the unstable
% eigenvalue.  Section 7 of the paper (sec:parametrix).
% Owned shared labels: def:transfer:parameter-choice, def:transfer:structural-data.
% Paper labels: prop:global-spectral-transfer (Proposition 7.1), lem:parametrix-error
% (Lemma 7.2), lem:global-contour-resolvent (Lemma 7.3), lem:global-riesz-nonzero (Lemma 7.4),
% rem:flexible-buffer-scales (Remark 7.5, no node; see the end of this chapter's first section).

This chapter transfers the instability from the selected cell to the torus.  The
parametrix has \emph{two} summands: the background resolvent of the
previous chapter carries everything away from the selected cell, and the rescaled seed
resolvent carries the cell itself.  The error is a sum of two diffusion commutators, both
supported in annuli of width $\asymp R_n$ where the horizontal velocity vanishes, so it is
$O(\sqrt\eps/R_n+\eps/R_n^2)=O(\sqrt{\ell_n}+\ell_n)$ and is made small by one last
enlargement of the base index.  The fixed test vector $f_*$ of the seed, rescaled and
placed, then separates the global Riesz projection from zero.

\section{The ordered choice of parameters}

\begin{definition}[The global parameter chain]
  \label{def:transfer:parameter-choice}
  \lean{FastDynamo.GlobalSeedInterface,
    FastDynamo.GlobalSeedInterface.axialShift,
    FastDynamo.GlobalSeedInterface.errorConstant,
    FastDynamo.GlobalSeedInterface.approximateInverseConstant,
    FastDynamo.GlobalSeedInterface.errorTarget,
    FastDynamo.GlobalSeedInterface.mul_errorTarget_div_lt_quarter,
    FastDynamo.GlobalParameterChoice,
    FastDynamo.GlobalParameterChoice.velocity,
    FastDynamo.GlobalParameterChoice.jointPlacement,
    FastDynamo.GlobalParameterChoice.eps0,
    FastDynamo.GlobalParameterChoice.eps0_formula,
    FastDynamo.GlobalParameterChoice.matchingCell,
    FastDynamo.GlobalParameterChoice.axialAmplitude_match,
    FastDynamo.GlobalParameterChoice.errorSmall_cell,
    FastDynamo.exists_globalParameterChoice,
    FastDynamo.nonempty_globalParameterChoice,
    FastDynamo.LocalSeedC2.toGlobalSeedInterface,
    FastDynamo.HorizontalMode.globalSeedInterface,
    FastDynamo.HorizontalMode.mainParameterChoice,
    FastDynamo.HorizontalMode.mainParameterChoice_eps0_formula,
    FastDynamo.HorizontalMode.kStar_le_mainParameterChoice_baseIndex}
  \leanok
  \uses{def:local:seed, prop:local-seed, def:velocity:scale-family, def:velocity:placement}
  The parameters are fixed in the order
  \[
    (v,I,\Gamma,\cG,g_*,f_*,L_*,q,\vartheta)\ \longrightarrow\ \tau
    \ \longrightarrow\ K_0 ,
  \]
  each stage depending only on the previous ones.  (The paper states this order in prose,
  in Section~3.1 and in the paragraph after the proof of Proposition~7.1; it does not display
  it.)  The seed
  data are Proposition~\ref{prop:local-seed}; the single geometric ratio satisfies
  $1<q<\sqrt{\delta_+/\delta_-}$ and $\vartheta$ is the fixed axial cut-off; the axial shift
  is fixed by \eqref{eq:background:axial-separation-choice}; and the base index $K_0$ is
  enlarged once, to satisfy simultaneously the five conditions: the ratio bound
  $q+K_0^{-1}<\sqrt{\delta_+/\delta_-}$, the chart packing, the core containment
  $R_v\ell_n<R_n/10$ (for which the paper gives the sufficient threshold
  $K_0>(10R_v)^2$), \eqref{eq:background:axial-separation-choice}'s diffusion clause and
  \eqref{eq:transfer:error-small}.  One final $K_0$ fixes the velocity $u$ and the threshold
  $\eps_0=\delta_+k_*^2/K_0^2$ ($=\delta_+/K_0^2$ in the paper, where $k_*=1$); only the matching index $n=n(\eps)$ depends on $\eps$
  afterwards.  The witness-support bound \eqref{eq:transfer:scaled-witness-support} is
  \emph{not} a sixth threshold: it follows from the core containment because $R_v\geq L_*$
  (Lemma~\ref{lem:transfer:witness}).
\end{definition}

In Lean the chain is a record \texttt{GlobalParameterChoice S errorConstant errorTarget}
over a \texttt{GlobalSeedInterface} $S$, and
\texttt{exists\_globalParameterChoice} inhabits it for every seed interface $S$, every
$\texttt{errorConstant}\geq0$ and every $\texttt{errorTarget}>0$: the
$K_0$-thresholds are collected into a single monotone funnel of the shape ``$\ell_0<t$ or
$\sqrt{\ell_0}<t$'' and discharged at once.  The record has exactly four fields --
\texttt{ratios} (the single ratio $q$), \texttt{jointPlacement} (the chart packing and the
core containment, both read at $R_v=\max(1+h,L_*)$), \texttt{diffusionSmall} and
\texttt{errorSmall} -- matching the paper's five conditions, since the two placement
conditions travel together and the integer-ratio bound is a field of the
\texttt{ScaleFamily} inside the placement.  The axial shift
$\tau=\texttt{axialShift}$ is a closed term in $\tau$-free and $K_0$-free data, which is
what makes the order non-circular.  Note that the
interface record itself stores $L_*$ and $L$ as bare reals and carries no seed facts, so
\texttt{nonempty\_globalParameterChoice} on its own is a \emph{shape} statement; the
chain is instantiated only at \texttt{LocalSeedC2.toGlobalSeedInterface}, and every seed
fact the transfer consumes ($f_*$, its norm, its support, the resolvents, the projection
bound) is taken from the \texttt{LocalSeedC2} record, not from the interface.  The chosen
chain is \texttt{HorizontalMode.mainParameterChoice}.

One point about the shape of the dependency graph here.  The record
\texttt{GlobalParameterChoice S errorConstant errorTarget} takes the Section~6 and
Section~7 constants
as \emph{arguments}: $\tau$ is a closed term in the seed data and the ratio, and the
thresholds \eqref{eq:background:axial-separation-choice} and \eqref{eq:transfer:error-small}
enter as fields whose numerical content the caller supplies.  That is exactly the paper's
quantifier order, and it is why this node depends
only on the seed and on the placement in the graph, while the \emph{instantiation}
\texttt{mainParameterChoice} additionally feeds it
\texttt{GlobalSeedInterface.errorConstant} and \texttt{errorTarget}, closed terms built from
the constants of Proposition~\ref{prop:background-resolvent},
Corollary~\ref{cor:scaled-local-resolvents} and Lemma~\ref{lem:parametrix-error}.  Drawing
those two extra edges would make the graph cyclic; the paper avoids the circle in the
same way, by fixing the constants before the base index.

Remark~7.5 of the paper (\texttt{rem:flexible-buffer-scales}), which observes that the
arguments of Sections~6 and~7 also work for other non-increasing buffer radii $R_n$ with
$\sum_nR_n\to0$ and $\sup_n\ell_n/R_n\to0$ as $K_0\to\infty$, has no node and is not
formalised: the Lean fixes $R_n=\sqrt{\ell_n}$ (\texttt{ScaleFamily.bufferRadius}), the
paper's own choice.

\section{The two-resolvent approximate inverse}

\begin{definition}[The three cut-offs and the chart transports]
  \label{def:transfer:structural-data}
  \lean{FastDynamo.ParametrixCutoffs.backgroundCutoff,
    FastDynamo.ParametrixExplicitCutoff.inputCutoff,
    FastDynamo.ParametrixExplicitCutoff.outputCutoff,
    FastDynamo.ParametrixCutoffs.norm_cutoff_coordinateDerivative_sum_le,
    FastDynamo.ParametrixCutoffs.norm_cutoff_iteratedFDeriv_two_sum_le,
    FastDynamo.ParametrixCutoffs.backgroundCutoff_mul_one_sub_inputCutoff,
    FastDynamo.ParametrixCutoffs.outputCutoff_mul_inputCutoff,
    FastDynamo.ParametrixCutoffs.cutoff_partition,
    FastDynamo.ParametrixCutoffs.TransverseVecSecondOrderCutoff,
    FastDynamo.ParametrixCutoffs.HorizontalParametrixCutoffs,
    FastDynamo.ParametrixCutoffs.HorizontalParametrixCutoffs.partition,
    FastDynamo.TorusPlacedCutoffC2.placedCutoffs,
    FastDynamo.ScaleFamily.TorusCellPlacement.placedParametrixCutoffs,
    FastDynamo.TorusChartC2Transfer.restrictInputC2,
    FastDynamo.TorusChartC2Transfer.extendStateC2,
    FastDynamo.TorusChartC2Transfer.extendDomainC2,
    FastDynamo.TorusChartC2Transfer.extendStateC2_comp_restrictInputC2}
  \leanok
  \uses{def:velocity:placement, def:transfer:parameter-choice}
  Choose smooth scalar cut-offs with values in $[0,1]$, the \emph{input cut-off} $\phi_n$ and
  the \emph{output cut-offs} $\psi_n$, $\chi_n$,
  \[
    \phi_n=1\text{ on }D(p_n,R_n/5),\quad\supp\phi_n\subset D(p_n,R_n/4),
    \qquad
    \psi_n=1\text{ on }D(p_n,R_n/3),\quad\supp\psi_n\subset D(p_n,R_n/2),
  \]
  \[
    \chi_n=0\text{ on }D(p_n,R_n/10),\qquad\chi_n=1\text{ outside }D(p_n,R_n/6),
  \]
  obtained by rescaling fixed radial cut-offs, so that
  \begin{equation}
    \norm{D^j\psi_n}_\infty+\norm{D^j\chi_n}_\infty\leq CR_n^{-j},
    \qquad j=1,2,
    \label{eq:transfer:cutoff-bounds}
  \end{equation}
  and
  \begin{equation}
    \chi_n(1-\phi_n)=1-\phi_n,\qquad\psi_n\phi_n=\phi_n .
    \label{eq:transfer:cutoff-identities}
  \end{equation}
\end{definition}

Only the input cut-off $\phi_n$ splits the forcing, so only the two output cut-offs need
derivative bounds; the paper's \texttt{eq:parametrix-cutoff-bounds} accordingly carries no
summand $\norm{D^j\phi_n}_\infty$.  Lean matches
this: the consumed interface \texttt{HorizontalParametrixCutoffs} stores, for $\phi_n$, only
the two sup-norm bounds $\norm{\phi_n}_\infty\leq1$ and $\norm{1-\phi_n}_\infty\leq1$, and
the four derivative bounds only for $\psi_n$ and $\chi_n$.  The auxiliary lemmas
\texttt{norm\_cutoff\_coordinateDerivative\_sum\_le} and
\texttt{norm\_cutoff\_iteratedFDeriv\_two\_sum\_le} do bound all three cut-offs in a single
three-term sum -- that is how the explicit construction produces them -- but the $\phi_n$
summand is never read: the placed construction extracts the $\psi_n$ and $\chi_n$ halves,
which is exactly the paper's display.

In Lean a cut-off is packaged as a \texttt{TransverseVecSecondOrderCutoff}: the
multiplication operators by the value, by each first derivative and by the Laplacian, with
their $\Lp^\infty$ bounds and the identity that the value operator maps $\Hs^2$ into
$\Hs^2$.  \texttt{HorizontalParametrixCutoffs R} bundles the three of them together with
the operator form $\chi_n(1-\phi_n)+\psi_n\phi_n=\id$ of
\eqref{eq:transfer:cutoff-identities}.  The concrete placed triple is built on the plane
and transported to the torus through \texttt{placedCutoffs}, which needs the no-wrap
statement that the closed ball of radius $R_n/2$ around $p_n$ lies in one fundamental box;
the chart transports $\texttt{restrictInputC2}$ (torus $\to$ plane, multiply by $\phi_n$,
extend by zero) and $\texttt{extendStateC2}$ (plane $\to$ torus, multiply by $\psi_n$) are
the two maps that make the selected summand well defined.

\begin{definition}[The approximate inverse]
  \label{def:transfer:approximate-inverse}
  \lean{FastDynamo.TorusTwoSummandParametrix.twoSummandParametrix,
    FastDynamo.TorusTwoSummandParametrix.twoSummandError,
    FastDynamo.TorusTwoSummandParametrix.BackgroundParametrixData,
    FastDynamo.TorusTwoSummandParametrix.BackgroundParametrixData.ofTorusBackgroundResolventData,
    FastDynamo.TorusTwoSummandParametrix.SelectedLocalResolventData,
    FastDynamo.TorusTwoSummandParametrix.TwoSummandParametrixData,
    FastDynamo.TorusTwoSummandParametrix.TwoSummandParametrixData.approximateInverse,
    FastDynamo.TorusTwoSummandParametrix.TwoSummandParametrixData.continuousOn_approximateInverse,
    FastDynamo.TorusChartC2Transfer.selectedLocalResolventData,
    FastDynamo.TorusChartC2Transfer.selectedCellParametrixData,
    FastDynamo.TorusParametrixError.norm_stateParametrix_apply_le}
  \leanok
  \uses{def:transfer:structural-data, prop:background-resolvent,
    cor:scaled-local-resolvents}
  With $\cR_{\rm bg}(z)=(z-A^{\rm bg}_{\eps,n})^{-1}$ and
  $\cR_n(z)=(z-A^{(n)}_{\eps,n})^{-1}$ -- the first an abbreviation the paper introduces in
  running text just before \texttt{eq:Rapp}, the second the resolvent of Corollary~6.2 of
  the paper -- put
  \begin{equation}
    \cR_{\rm app}(z)=\chi_n\cR_{\rm bg}(z)(1-\phi_n)+\psi_n\cR_n(z)\phi_n,
    \qquad z\in\Gamma_\tau .
    \label{eq:transfer:Rapp}
  \end{equation}
  For the second summand, $\phi_nf$ is extended by zero from the fixed chart to $\R^2$, the
  whole-plane resolvent is applied, and the product with $\psi_n$ is restricted to the
  torus.  Both summands map $\Lp^2(\T^2;\CC^2)$ boundedly into $\Hs^2(\T^2;\CC^2)$, and
  \eqref{eq:background:bg-resolvent}, \eqref{eq:background:local-gradient} give
  \begin{equation}
    \sup_{z\in\Gamma_\tau}\norm{\cR_{\rm app}(z)}_{2\to2}\leq C_{\rm app},
    \label{eq:transfer:parametrix-uniform-bound}
  \end{equation}
  uniformly in $n$, in $\delta_n\in I$ and in the final enlargement of $K_0$.  The
  $\Lp^2\to\Hs^2$ norm need not be uniform.
\end{definition}

The Lean \texttt{TwoSummandParametrixData} is an interface: the two resolvent data, the
cut-off triple, the chart transports and the hypotheses that make the commutator
computation valid.  Its instance at the selected cell is
\texttt{selectedCellParametrixData}.  The uniform bound is assembled as
$C_{\rm app}=C_{\rm bg}+C_{\rm loc}$
(\texttt{TorusParametrixError.norm\_stateParametrix\_apply\_le}), using the sharp
$\norm{1-\phi_n}_{\Lp^\infty}\leq1$, which holds because $0\leq\phi_n\leq1$.

\begin{lemma}[The commutator of the operator with a cut-off]
  \label{lem:transfer:commutator}
  \lean{FastDynamo.HorizontalCutoffCommutator.commutator,
    FastDynamo.HorizontalCutoffCommutator.diffusion_commutator_apply,
    FastDynamo.HorizontalCutoffCommutator.axialTransport_commutes,
    FastDynamo.HorizontalCutoffCommutator.stretching_commutes,
    FastDynamo.HorizontalCutoffCommutator.commutator_apply,
    FastDynamo.HorizontalCutoffCommutator.VelocityVanishesOnGradient,
    FastDynamo.HorizontalCutoffCommutator.commutator_apply_of_velocityVanishesOnGradient,
    FastDynamo.HorizontalCutoffCommutator.norm_commutator_apply_le_explicitScale}
  \leanok
  \uses{def:operator:horizontal-induction, def:transfer:structural-data}
  Write $[A_{\eps,K_n},\chi]=A_{\eps,K_n}\chi-\chi A_{\eps,K_n}$.  For a scalar cut-off $\chi$
  whose gradient is supported where $u_\perp$ vanishes,
  \[
    [A_{\eps,K_n},\chi]h
      =\eps\bigl(2\nabla_\perp\chi\cdot\nabla_\perp h+(\Delta_\perp\chi)h\bigr),
    \qquad
    \norm{[A_{\eps,K_n},\chi]h}_2\leq C\Bigl(\frac\eps{R_n}\norm{\nabla_\perp h}_2
      +\frac\eps{R_n^2}\norm h_2\Bigr).
  \]
\end{lemma}

\begin{proof}
  \leanok
  The axial multiplier and the stretching term are of order zero and commute with
  multiplication by the scalar $\chi$.  The transport term contributes
  $(u_\perp\cdot\nabla_\perp\chi)h$, which vanishes because $u_\perp=0$ on the support of
  $\nabla_\perp\chi$.  The commutator is therefore that of $\eps\Delta_\perp$, which the
  Leibniz rule turns into the displayed identity, and the bound follows from
  $\norm{\nabla_\perp\chi}_{\Lp^\infty}\leq C/R_n$ and
  $\norm{\Delta_\perp\chi}_{\Lp^\infty}\leq C/R_n^2$.
\end{proof}

The axial multiplier and the stretching term are zeroth order, so they commute with a
scalar cut-off \emph{exactly}; this is proved in Lean rather than asserted.  The transport
term is killed by the hypothesis \texttt{VelocityVanishesOnGradient}, which is the only
geometric input and is discharged for $\chi_n$ (gradient in $R_n/10\leq\abs{y-p_n}\leq
R_n/6$) and for $\psi_n$ (gradient in $R_n/3\leq\abs{y-p_n}\leq R_n/2$) from the support
fit $\supp u_n^{\rm loc}\subset D(p_n,R_n/10)$, in one statement covering both.  The paper
proves this commutator identity and bound inside the proof of Lemma~7.2, for
$\chi\in\{\chi_n,\psi_n\}$.

\begin{lemma}[Identity and error for the approximate inverse]
  \label{lem:parametrix-error}
  \lean{FastDynamo.TorusTwoSummandParametrix.TwoSummandParametrixData.right_error_identity,
    FastDynamo.TorusTwoSummandParametrix.TwoSummandParametrixData.error_eq,
    FastDynamo.TorusTwoSummandParametrix.TwoSummandParametrixData.transfer_rightError_eq,
    FastDynamo.TorusParametrixError.scaleErrorMajorant,
    FastDynamo.TorusParametrixError.commutatorErrorConstant,
    FastDynamo.TorusParametrixError.parametrixErrorConstant,
    FastDynamo.TorusParametrixError.norm_error_apply_le,
    FastDynamo.TorusParametrixError.norm_error_le,
    FastDynamo.TorusParametrixError.scaleErrorMajorant_scaled_le}
  \leanok
  \uses{def:transfer:approximate-inverse, lem:transfer:commutator,
    prop:background-resolvent, cor:scaled-local-resolvents}
  For $z\in\Gamma_\tau$,
  \begin{equation}
    (z-A_{\eps,K_n})\cR_{\rm app}(z)=\id+E_n(z),
    \qquad
    E_n(z)=-[A_{\eps,K_n},\chi_n]\cR_{\rm bg}(z)(1-\phi_n)
      -[A_{\eps,K_n},\psi_n]\cR_n(z)\phi_n,
    \label{eq:transfer:parametrix}
  \end{equation}
  and
  \begin{equation}
    \norm{E_n(z)}_{2\to2}\leq C\Bigl(\frac{\sqrt\eps}{R_n}+\frac\eps{R_n^2}\Bigr)
      \leq C(\sqrt{\ell_n}+\ell_n),
    \label{eq:transfer:explicit-error}
  \end{equation}
  with $C$ independent of $n$, of $\delta_n\in I$, of $z$ and of the final enlargement of
  $K_0$.
\end{lemma}

The paper gives the formula for $E_n(z)$ in the proof; the statement above includes it.
The Lean constant is explicit: $C=(2\sqrt2C_1+2C_2)(C_{\rm bg}+C_{\rm loc})$, with $C_1$,
$C_2$ the first- and second-derivative constants of the rescaled cut-offs, and the scaling
identity $\eps=\delta\ell^2$, $R=\sqrt\ell$ turns $C(\sqrt\eps/R+\eps/R^2)$ into
$C(1+\delta_+)(\sqrt\ell+\ell)$.

\begin{proof}
  \leanok
  \uses{lem:transfer:commutator, def:transfer:approximate-inverse}
  The full and background coefficients agree wherever $\chi_n\neq0$, because
  $\supp u_n^{\rm loc}\subset D(p_n,R_n/10)$; the full and selected whole-plane
  coefficients agree on $\supp\psi_n$, because the added axial velocity is constant on
  $D(p_n,3R_n/4)$ and all other cells are disjoint from the buffer disc.  So
  \eqref{eq:transfer:cutoff-identities} gives
  $(z-A_{\eps,K_n})\chi_n\cR_{\rm bg}(1-\phi_n)
  =1-\phi_n-[A_{\eps,K_n},\chi_n]\cR_{\rm bg}(1-\phi_n)$ and
  $(z-A_{\eps,K_n})\psi_n\cR_n\phi_n=\phi_n-[A_{\eps,K_n},\psi_n]\cR_n\phi_n$; adding gives
  \eqref{eq:transfer:parametrix}.  Lemma~\ref{lem:transfer:commutator} bounds each
  commutator by $C(\eps/R_n)\norm{\nabla_\perp h}_2+C(\eps/R_n^2)\norm h_2$, and the two
  resolvent estimates \eqref{eq:background:bg-resolvent},
  \eqref{eq:background:local-gradient} convert
  $\eps\norm{\nabla_\perp h}_2\leq\sqrt\eps\cdot C\norm f_2$ and
  $\norm h_2\leq C\norm f_2$.
\end{proof}

\section{The global resolvent on the contour}

\begin{lemma}[The resolvent of the horizontal operator on $\T^2$ along the shifted contour]
  \label{lem:global-contour-resolvent}
  \lean{FastDynamo.TorusParametrixError.error_norm_lt_one_of_majorant,
    FastDynamo.TorusParametrixError.exists_scale_threshold,
    FastDynamo.TorusParametrixError.globalContourResolvent,
    FastDynamo.TorusParametrixError.globalContourResolvent_inv,
    FastDynamo.TorusChartC2Transfer.selectedCellGlobalContourResolvent,
    FastDynamo.TorusGlobalRieszNonzero.cellContourResolvent,
    FastDynamo.HorizontalMode.mainSelectedCellResolvent,
    FastDynamo.GlobalParameterChoice.scaleErrorMajorant_lt_errorTarget}
  \leanok
  \uses{lem:parametrix-error, lem:torus-sector-basic, def:transfer:parameter-choice}
  For all sufficiently large $K_0$, uniformly in $n\geq0$ and in
  $\eps=\delta_n\ell_n^2$ with $\delta_n\in I$, one has
  $\Gamma_\tau\subset\rho(A_{\eps,K_n})$ and
  \begin{equation}
    (z-A_{\eps,K_n})^{-1}=\cR_{\rm app}(z)\bigl(\id+E_n(z)\bigr)^{-1},
    \qquad z\in\Gamma_\tau .
    \label{eq:transfer:global-resolvent}
  \end{equation}
  The threshold is: choose $0<\eta_*<1/2$ with
  $\tfrac{\operatorname{length}(\Gamma)}{2\pi}C_{\rm app}\tfrac{\eta_*}{1-\eta_*}<\tfrac14$
  and enlarge $K_0$ until
  \begin{equation}
    \sup_{n\geq0}\sup_{\delta\in I}\sup_{z\in\Gamma_\tau}\norm{E_{n,\delta}(z)}_{2\to2}
      <\eta_* .
    \label{eq:transfer:error-small}
  \end{equation}
\end{lemma}

The paper chooses $\eta_*$ and states \texttt{eq:error-small} in the proof of Lemma~7.3; the
statement above includes them.  The strict inequality of \eqref{eq:transfer:error-small}
is the strict inequality in Lean as well.  The Lean threshold is
$\eta_*=\min\bigl(\tfrac12,(100(1+K))^{-1}\bigr)$ with $K=(1+r_\Gamma)(1+C_{\rm app})$
(\texttt{GlobalSeedInterface.errorTarget}), which gives $A\eta_*/(1-\eta_*)<\tfrac14$ for
every prefactor $0\leq A\leq10K$ (\texttt{mul\_errorTarget\_div\_lt\_quarter}), in
particular for $A=r_\Gamma C_{\rm app}$; the prefactor
$\operatorname{length}(\Gamma)/2\pi$ is exactly the radius $r_\Gamma$.

\begin{proof}
  \leanok
  \uses{lem:parametrix-error, lem:torus-sector-basic}
  Since $\ell_n\leq\ell_0=k_*/K_0$ ($=1/K_0$ in the paper, where $k_*=1$),
  \eqref{eq:transfer:explicit-error} tends to zero
  uniformly as $K_0\to\infty$, so \eqref{eq:transfer:error-small} can be arranged; the
  Neumann series then inverts $\id+E_n(z)$ and \eqref{eq:transfer:parametrix} exhibits
  $\cR_{\rm app}(\id+E_n)^{-1}$ as a right inverse with values in $\Hs^2$.  By
  Lemma~\ref{lem:torus-sector-basic} the pencil $z-A_{\eps,K_n}\colon\Hs^2\to\Lp^2$ is
  Fredholm of index zero, so surjectivity forces injectivity and the right inverse is the
  resolvent; continuity in $z$ is inherited from the two resolvent families and the
  Neumann series.
\end{proof}

\section{The Riesz projection of the horizontal operator on $\T^2$ is non-zero}

\begin{lemma}[The placed witness]
  \label{lem:transfer:witness}
  \lean{FastDynamo.TorusSelectedWitness.VanishesOutsideDisk,
    FastDynamo.TorusSelectedWitness.scaledWitness,
    FastDynamo.TorusSelectedWitness.placedWitness,
    FastDynamo.TorusSelectedWitness.norm_placedWitness,
    FastDynamo.TorusSelectedWitness.inputMul_stateMul_placedWitness,
    FastDynamo.TorusSelectedWitness.outputMul_stateMul_placedWitness,
    FastDynamo.TorusSelectedWitness.restrictInputC2_placedWitness,
    FastDynamo.TorusSelectedWitness.vanishesOutsideDisk_scaledWitness,
    FastDynamo.GlobalParameterChoice.witnessSupport}
  \leanok
  \uses{prop:local-seed, def:transfer:structural-data, def:transfer:parameter-choice}
  Since $L_*\leq R_v$, the core containment $R_v\ell_n<R_n/10$ of
  Definition~\ref{def:velocity:placement} gives
  \begin{equation}
    L_*\ell_n<R_n/5\qquad\text{for every }n\geq0 .
    \label{eq:transfer:scaled-witness-support}
  \end{equation}
  Read $f_n=U_nf_*$ as a torus field by restriction to the fixed chart and extension by
  zero.  Then $\norm{f_n}_2=1$ and $\phi_nf_n=\psi_nf_n=f_n$.
\end{lemma}

\begin{proof}
  \leanok
  The placement radius satisfies $R_v\geq L_*$ by construction
  (\texttt{GlobalSeedInterface.witnessSupportRadius\_le\_placementRadius}), so
  $L_*\ell_n\leq R_v\ell_n<R_n/10<R_n/5$; the support of $f_*$ lies in $D(0,L_*)$, so the
  support of the rescaled and placed $f_n$ lies in $D(p_n,L_*\ell_n)\subset D(p_n,R_n/5)$,
  where both $\phi_n$ and $\psi_n$ equal one, and unitarity of $U_n$ gives $\norm{f_n}_2=1$.
\end{proof}

The radius $R_n/5$ in \eqref{eq:transfer:scaled-witness-support} is exactly the plateau
radius of the input cut-off, so the constant is not slack.  The paper requires
$R_v\geq L_*$ when $R_v$ is chosen (Section~3.3, just before \texttt{eq:local-support-in-buffer})
and derives the bound, its \texttt{eq:scaled-witness-support}, from the core containment
at the start of the proof of Lemma~7.4 rather than from a further enlargement of $K_0$,
which keeps the list of $K_0$-thresholds at five.
Lean does the same: the placement is chosen once, at support radius
$R_v=\max(1+h,L_*)$, and \texttt{GlobalParameterChoice.witnessSupport} is proved
from the placement's support fit together with $R_n/10<R_n/5$.

\begin{lemma}[The contour identity at the witness]
  \label{lem:transfer:contour-identity}
  \lean{FastDynamo.TorusSelectedContourIdentity.SelectedCellData,
    FastDynamo.TorusSelectedContourIdentity.SelectedCellData.parametrix,
    FastDynamo.TorusSelectedContourIdentity.circleIntegral_translate,
    FastDynamo.TorusSelectedContourIdentity.inner_stateParametrix_eq_of_witness,
    FastDynamo.TorusSelectedContourIdentity.inner_stateParametrix_placedWitness,
    FastDynamo.TorusSelectedContourIdentity.inner_contourOperator_placedWitness,
    FastDynamo.TorusGlobalRieszNonzero.re_inner_contourOperator_placedWitness_gt_half}
  \leanok
  \uses{lem:transfer:witness, def:transfer:approximate-inverse, lem:exact-renormalization,
    prop:local-seed}
  For $z\in\Gamma_\tau$,
  $\inner{\cR_{\rm app}(z)f_n}{f_n}=\inner{(z+i\tau-\cL_{\delta_n})^{-1}f_*}{f_*}$, and
  therefore
  \begin{equation}
    \frac1{2\pi i}\int_{\Gamma_\tau}\inner{\cR_{\rm app}(z)f_n}{f_n}\dd z
      =\inner{P_{\delta_n}f_*}{f_*} .
    \label{eq:transfer:selected-contour-identity}
  \end{equation}
  In particular $\re$ of the left-hand side exceeds $\tfrac12$.
\end{lemma}

\begin{proof}
  \leanok
  \uses{lem:transfer:witness, lem:exact-renormalization, prop:local-seed}
  The input of the background summand is $(1-\phi_n)f_n=0$, and the output cut-off of the
  selected summand equals one on $\supp f_n$, so only the selected summand survives; the
  scaling identity \eqref{eq:operator:exact-scaling} and the unitarity of $U_n$ turn it into
  the seed resolvent at $z+i\tau$.  Integrating and translating the circle by $+i\tau$ --
  the parametrisation and its derivative are unchanged, so the translation is exact and
  orientation preserving -- gives the seed's Riesz projection, and
  \eqref{eq:local:spectral-witness} gives the bound.
\end{proof}

\begin{lemma}[The Riesz projection of the horizontal operator on $\T^2$ is non-zero]
  \label{lem:global-riesz-nonzero}
  \lean{FastDynamo.TorusGlobalRieszNonzero.norm_stateResolvent_sub_stateParametrix_le_of_inv,
    FastDynamo.TorusGlobalRieszNonzero.norm_rieszProjection_sub_contourOperator_le_cell,
    FastDynamo.TorusGlobalRieszNonzero.re_inner_rieszProjection_placedWitness_gt_quarter,
    FastDynamo.TorusGlobalRieszNonzero.rieszProjection_ne_zero_cell,
    FastDynamo.LocalSeedC2.projectionError_matchingCellData_lt_quarter}
  \leanok
  \uses{lem:transfer:contour-identity, lem:global-contour-resolvent,
    lem:parametrix-error}
  The projection
  $Q_{\eps,n}=(2\pi i)^{-1}\int_{\Gamma_\tau}(z-A_{\eps,K_n})^{-1}\dd z$ satisfies
  \begin{equation}
    \Bigl|\inner{Q_{\eps,n}f_n}{f_n}-\inner{P_{\delta_n}f_*}{f_*}\Bigr|
      \leq\frac{\operatorname{length}(\Gamma)}{2\pi}C_{\rm app}
        \frac{\eta_*}{1-\eta_*}<\frac14,
    \label{eq:transfer:global-projection-error}
  \end{equation}
  hence $\re\inner{Q_{\eps,n}f_n}{f_n}>\tfrac14$ and $Q_{\eps,n}\neq0$, for every $n\geq0$
  and every $\eps=\delta_n\ell_n^2$ with $\delta_n\in I$.
\end{lemma}

The paper's statement of Lemma~7.4 asserts only $Q_{\eps,n}\neq0$; the bound
\eqref{eq:transfer:global-projection-error} (the paper's \texttt{eq:global-projection-error})
and $\re\inner{Q_{\eps,n}f_n}{f_n}>\tfrac14$ are established in its proof, together with
the content of Lemmas~\ref{lem:transfer:witness} and~\ref{lem:transfer:contour-identity}.
In Lean the error is bounded in operator norm,
$\norm{Q_{\eps,n}-(2\pi i)^{-1}\int_{\Gamma_\tau}\cR_{\rm app}(z)\dd z}
\leq r_\Gamma C_{\rm app}e/(1-e)$ with $e<1$ the explicit majorant of
$\sup_{z}\norm{E_n(z)}$ (\texttt{norm\_rieszProjection\_sub\_contourOperator\_le\_cell}),
then $e<\eta_*$ is used and the result is paired with the unit vector $f_n$.

\begin{proof}
  \leanok
  \uses{lem:global-contour-resolvent, lem:transfer:contour-identity}
  From \eqref{eq:transfer:global-resolvent},
  $(z-A_{\eps,K_n})^{-1}-\cR_{\rm app}(z)=-\cR_{\rm app}(z)(\id+E_n(z))^{-1}E_n(z)$, whose
  norm is at most $C_{\rm app}\eta_*/(1-\eta_*)$ by
  \eqref{eq:transfer:parametrix-uniform-bound} and
  \eqref{eq:transfer:error-small}; integrating over a circle of radius $r_\Gamma$ and
  pairing with the unit vector $f_n$ gives
  \eqref{eq:transfer:global-projection-error}.  Subtracting from
  Lemma~\ref{lem:transfer:contour-identity} leaves
  $\re\inner{Q_{\eps,n}f_n}{f_n}>\tfrac12-\tfrac14$.
\end{proof}

\section{The spectral transfer}

\begin{lemma}[Extraction of an eigenpair from a non-zero Riesz projection]
  \label{lem:transfer:eigenpair-extraction}
  \lean{FastDynamo.HorizontalSpectralExtraction.isCompactOperator_rieszProjection_horizontalInduction,
    FastDynamo.HorizontalSpectralExtraction.finiteDimensional_rieszProjection_range_horizontalInduction,
    FastDynamo.HorizontalSpectralExtraction.exists_horizontalInduction_eigenpair_mem_closedRegion_of_rieszProjection_ne_zero,
    FastDynamo.HorizontalSpectralExtraction.exists_growthFloor_horizontalInduction_eigenpair_of_rieszProjection_ne_zero,
    FastDynamo.ParametrixNeumannCorrection.ne_zero_of_inner_contour_test}
  \leanok
  \uses{lem:torus-sector-basic, lem:global-riesz-nonzero,
    lem:operator:shifted-contour}
  If $A_{\eps,K}$ has a resolvent on a spectral circle and the associated Riesz projection
  is non-zero, then $A_{\eps,K}$ has an eigenvalue $\lambda$ in the closed interior of that
  circle, with an eigenfunction in $\Hs^2(\T^2;\CC^2)$; in particular
  $\re\lambda\geq$ the circle's growth floor.
\end{lemma}

\begin{proof}
  \leanok
  \uses{lem:torus-sector-basic}
  The state resolvent is compact at every resolvent point
  (Lemma~\ref{lem:torus-sector-basic}), so $Q$ is a compact idempotent and has
  finite-dimensional range; a non-zero finite-rank idempotent built from a contour integral
  forces a spectral point inside, and index zero turns that spectral point into an
  eigenvalue with an $\Hs^2$ eigenfunction.  The real-part bound is the definition of the
  growth floor, which is invariant under the translation by $-i\tau$.
\end{proof}

\begin{proposition}[Existence of an unstable eigenvalue of the horizontal operator on $\T^2$]
  \label{prop:global-spectral-transfer}
  \lean{FastDynamo.TorusGlobalSpectralTransfer.exists_growing_horizontalInduction_eigenpair_cell,
    FastDynamo.TorusGlobalSpectralTransfer.exists_growing_reducedInduction_eigenpair_cell,
    FastDynamo.LocalSeedC2.matchingCellData,
    FastDynamo.LocalSeedC2.exists_growing_reducedInduction_eigenpair_matchingCell,
    FastDynamo.HorizontalMode.exists_growing_reducedInduction_eigenpair_matchingCell,
    FastDynamo.mainReducedClauses_holds}
  \leanok
  \uses{def:transfer:parameter-choice, lem:global-contour-resolvent,
    lem:global-riesz-nonzero, lem:transfer:eigenpair-extraction, lem:axial-recovery}
  Fix $v,I,\Gamma,\cG,g_*,f_*,L_*$, and choose $q$, the axial cut-off $\vartheta$ and
  $\tau$ in the order of Definition~\ref{def:transfer:parameter-choice}.
  Then $K_0$ can be chosen so that for
  every $0<\eps\leq\eps_0$ and every $n$ with $\delta_n=\eps/\ell_n^2\in I$,
  \[
    \Gamma_\tau\subset\rho(A_{\eps,K_n}),\qquad
    Q_{\eps,n}=\frac1{2\pi i}\int_{\Gamma_\tau}(z-A_{\eps,K_n})^{-1}\dd z\neq0 .
  \]
  In particular $A_{\eps,K_n}$ has an eigenvalue inside $\Gamma_\tau$ with real part at least
  $g_*$:
  \[
    \lambda^\eps=\alpha_\eps+i\omega_\eps,\qquad\alpha_\eps\geq g_* .
  \]
\end{proposition}

The paper's statement is existential (``has an eigenvalue inside $\Gamma_\tau$ with real
part at least $g_*$''), and so is the Lean,
\texttt{exists\_growthFloor\_horizontalInduction\_eigenpair\_of\_rieszProjection\_ne\_zero};
the paper does not name the eigenvalue, and the display above only names it for later use.
The paper's proof also notes the universal form, ``every eigenvalue in this region has real
part at least $g_*$, as translation by $-i\tau$ preserves real parts''.  In Lean that is
one application of \texttt{SpectralCircle.growthFloor\_le\_re\_of\_mem\_ball} -- the very
lemma the existential proof applies -- and nothing downstream needs it, so it is not stated
separately.

In Lean the proposition is used in its \emph{eigenpair} form and immediately composed with
the axial recovery, so the two published endpoints are
\texttt{TorusGlobalSpectralTransfer.exists\_growing\_reducedInduction\_eigenpair\_cell}
(one selected cell, three components) and, at the chosen chain,
\texttt{HorizontalMode.exists\_growing\_reducedInduction\_eigenpair\_matchingCell}, whose
packaging over all $0<\eps\leq\eps_0$ is \texttt{mainReducedClauses\_holds}.  The last of
these is the exact content of this node at the constructed velocity; it is the theorem the
whole development exists to prove.

\begin{proof}
  \leanok
  \uses{lem:global-contour-resolvent, lem:global-riesz-nonzero,
    lem:transfer:eigenpair-extraction, lem:torus-sector-basic}
  Lemma~\ref{lem:global-contour-resolvent} gives the contour resolvent and
  Lemma~\ref{lem:global-riesz-nonzero} the non-zero projection, at every admissible cell
  and every admissible diffusivity.  By Lemma~\ref{lem:torus-sector-basic} the resolvent is
  compact, so the spectrum inside $\Gamma_\tau$ consists of isolated eigenvalues of finite
  algebraic multiplicity, and the non-zero projection shows that it is non-empty;
  Lemma~\ref{lem:transfer:eigenpair-extraction} extracts the eigenpair.  Every eigenvalue in
  the region has $\re\lambda\geq g_*$ because translation by $-i\tau$ preserves real
  parts.
\end{proof}

\begin{lemma}[Every admissible cell]
  \label{lem:transfer:all-cells}
  \lean{FastDynamo.exists_growing_eigenpair_of_matchingCell,
    FastDynamo.exists_growing_eigenpair_of_admissible_index,
    FastDynamo.MainEveryAdmissibleCellClauses,
    FastDynamo.mainEveryAdmissibleCellClauses_holds}
  \leanok
  \uses{prop:global-spectral-transfer, def:transfer:parameter-choice}
  Proposition~\ref{prop:global-spectral-transfer} holds not only at a canonically chosen
  matching cell but at every cell index $n$ with $\eps/\ell_n^2\in I$: the axial index
  $K_n$ is non-zero and carries an eigenpair of $A_{\eps,K_n}$, hence of the
  three-component operator, with $\re\lambda\geq\gamma_0$.
\end{lemma}

This is stronger than Theorem~\ref{thm:main} requires and is recorded as item (4) of
Theorem~\ref{thm:main-strengthened}.  The Lean predicate quantifies over
$0<\eps\leq\eps_0$, although the paper's proof needs no upper bound on $\eps$ once an
admissible index is given; the threshold enters only through the \emph{existence} of such
an index.

\begin{proof}
  \leanok
  \uses{prop:global-spectral-transfer}
  Every step of the chain is stated at an arbitrary admissible index, so nothing has to be
  redone; only the choice of the index is dropped.
\end{proof}
